\section{Fase 5, Verificación del Contenido de los Archivos}

\textbf{Objetivo:} Confirmar que los archivos \texttt{.iq} generados contienen datos válidos, metadatos correctos, y que incluyen tanto el pre-evento como el post-evento capturados en el buffer circular.

% --- TEST 11 ---
\subsection{T-11, Metadatos del cabecero correctos}

Ejecute el siguiente script sobre uno de los archivos guardados en T-07:

\begin{lstlisting}[language=bash]
python3 - <<'EOF'
import struct, json, sys
path = sorted(__import__('glob').glob(
    '/home/kraken/Signal-Recorder/events/event_Iridium_*.iq'))[0]
with open(path, 'rb') as f:
    hlen = struct.unpack('I', f.read(4))[0]
    h = json.loads(f.read(hlen))
    f.read(8)  # timestamp
    n = struct.unpack('I', f.read(4))[0]

print(f"File:          {path}")
print(f"Group:         {h['group_name']}")
print(f"Center freq:   {h['center_freq']/1e6:.4f} MHz")
print(f"Target freq:   {h['target_freq']/1e6:.4f} MHz")
print(f"Sample rate:   {h['sample_rate']/1e6:.2f} MHz")
print(f"Gain:          {h['gain']} dB")
print(f"Pre samples:   {h['pre_samples']:,}  ({h['pre_samples']/h['sample_rate']:.2f} s)")
print(f"Post samples:  {h['post_samples']:,}  ({h['post_samples']/h['sample_rate']:.2f} s)")
print(f"Peak power:    {h['peak_power_db']:.1f} dB")
print(f"Baseline:      {h['baseline_power_db']:.1f} dB")
print(f"Total samples: {n:,}")
EOF
\end{lstlisting}

\begin{tcolorbox}[colback=gray!5, colframe=deyposcolor, title=Resultado esperado (T-11)]
\begin{itemize}
    \item \texttt{Center freq} = \textbf{1621.3000\,MHz}
    \item \texttt{Target freq} = \textbf{1621.3000\,MHz}
    \item \texttt{Sample rate} = \textbf{1.50\,MHz}
    \item \texttt{Gain} = \textbf{40.0\,dB}
    \item \texttt{Pre samples} $\approx$ 4\,500\,000 ($\approx$ 3,0\,s), puede ser ligeramente inferior si el grabador arrancó recientemente
    \item \texttt{Post samples} $\approx$ 4\,500\,000 ($\approx$ 3,0\,s)
    \item \texttt{Peak power} notablemente superior a \texttt{Baseline} (diferencia $>$ \texttt{threshold\_db} configurado)
\end{itemize}
\end{tcolorbox}

% --- TEST 12 ---
\subsection{T-12, Datos IQ no nulos y finitos}

\begin{lstlisting}[language=bash]
python3 - <<'EOF'
import struct, json, numpy as np, glob
path = sorted(glob.glob(
    '/home/kraken/Signal-Recorder/events/event_Iridium_*.iq'))[0]
with open(path, 'rb') as f:
    hlen = struct.unpack('I', f.read(4))[0]
    f.read(hlen); f.read(8)
    n = struct.unpack('I', f.read(4))[0]
    raw = np.frombuffer(f.read(n * 2 * 4), dtype=np.float32)

samples = raw[0::2] + 1j * raw[1::2]
print(f"Total samples:  {len(samples):,}")
print(f"All finite:     {np.all(np.isfinite(samples))}")
print(f"Non-zero:       {np.any(samples != 0)}")
print(f"RMS amplitude:  {np.sqrt(np.mean(np.abs(samples)**2)):.6f}")
print(f"Peak amplitude: {np.max(np.abs(samples)):.6f}")
EOF
\end{lstlisting}

\begin{tcolorbox}[colback=gray!5, colframe=deyposcolor, title=Resultado esperado (T-12)]
\begin{itemize}
    \item \texttt{All finite: True}
    \item \texttt{Non-zero: True}
    \item \texttt{RMS amplitude} $>$ 0 (típicamente entre 0,001 y 0,3 según la ganancia y el nivel de señal)
\end{itemize}
\end{tcolorbox}

% --- TEST 13 ---
\subsection{T-13, Gráfico \textit{waterfall} renderiza sin errores}

\begin{lstlisting}[language=bash]
source venv/bin/activate
IQ_FILE=$(ls /home/kraken/Signal-Recorder/events/event_Iridium_*_ch0.iq | head -1)
python src/watfall_plot.py "$IQ_FILE"
\end{lstlisting}

\begin{tcolorbox}[colback=gray!5, colframe=deyposcolor, title=Resultado esperado (T-13)]
Ningún error de Python. Aparece un archivo PNG en el directorio \texttt{./plots/}. Al abrirlo (puede copiarlo al portátil con \texttt{scp}), debe ser visible una ráfaga de señal en torno a la frecuencia central en el gráfico \textit{waterfall}.
\end{tcolorbox}
